The closest vector problem with preprocessing (CVPP)~\cite{Micciancio2001} allows all
lattice-dependent work to be performed before the target is known.
We show that exact decision-CVPP on the fixed principal-ideal
family below has no polynomial-time solution unless $\NP\subseteq\Ppoly$.

For a fixed lattice family, the preprocessing string is common to all target
queries and can serve as a nonuniform advice string depending
only on the input length.  We therefore encode every X3C
instance of a given size using only the target and threshold.

\paragraph{The X3C encoding.}
For $m\geq3$, let $k_m:=\binom m3$, and let
$\mathbf U_m\in\{0,1\}^{m\times k_m}$ be the matrix whose columns are the
incidence vectors of all three-element subsets of $[m]$,
listed in lexicographic order.  A mask
$\boldsymbol\mu\in\{0,1\}^{k_m}$ specifies which of these triples belong to
the input collection.  The collection has an exact cover precisely when
\[
 \mathbf U_m\boldsymbol\xi=\mathbf 1_m
 \quad\text{for some}\quad
 \boldsymbol\xi\in\{0,1\}^{k_m}
 \text{ with }\boldsymbol\xi\leq\boldsymbol\mu.
\]
Thus $(\mathbf U_m,\mathbf 1_m)$ depends only on $m$.  The mask
$\boldsymbol\mu$ encodes the input collection.

\Needspace{14\baselineskip}
\begin{theorem}
\label{thm:fixed-family}
There exists a polynomial-time reduction from X3C to exact decision-CVP such
that every X3C instance on a universe of size $m$ is mapped to an instance
$(I_m,T_{\boldsymbol\mu},\Delta_{\boldsymbol\mu})$, where
\(I_m\mathrel{:=}(g_m)\lhd R_{d_m}\) is a full-rank principal ideal and
its multiplication basis \(\mathbf B_{g_m}\) depends only on \(m\).
\end{theorem}

\begin{proof}
Apply Theorem~\ref{thm:compiler} to the fixed system
$(\mathbf U_m,\mathbf 1_m)$.  The reduction produces the ideal
\(I_m\) and its multiplication basis independently of the mask.  The target
\(T_{\boldsymbol\mu}\) and threshold \(\Delta_{\boldsymbol\mu}\) encode the
collection.
In the notation of Section~\ref{sec:compiler}, set \(M:=m\)
and \(n:=k_m\).
Thus $k_m=\Theta(m^3)$, and the explicit Sidon construction gives
$c=\alpha_{k_m}=\Theta(k_m^2)$.  Since $\rho_m+c=\Theta(mc)$, rounding up
to a power of two gives
 $d_m=\Theta(mk_m^2)=\Theta(m^7)$.  The closest squared distance is
$\Delta_{\boldsymbol\mu}$ for a YES instance and at least
$\Delta_{\boldsymbol\mu}+4$ for a NO instance.
\end{proof}

Together with the certificate argument of
Theorem~\ref{thm:cyclotomic-cvp}, Theorem~\ref{thm:fixed-family} also shows
that exact decision-CVP on this fixed family is $\NP$-complete.  Its main consequence
is the following preprocessing lower bound.

\begin{corollary}
\label{cor:cvpp}
If exact decision-CVPP is solvable in polynomial time on the
family $\{\mathcal L(I_m)\}_{m\geq3}$, each of which is represented by the
corresponding multiplication basis from Theorem~\ref{thm:fixed-family}, then
\[
 \NP\subseteq\Ppoly.
\]
\end{corollary}

\begin{proof}
Let \((\mathsf{Pre},\mathsf{Dec},p)\) be an exact
decision-CVPP scheme for \(\{\mathcal L(I_m)\}_{m\geq3}\), with each lattice
represented by its multiplication basis \(\mathbf B_{g_m}\).
Define the language
\(L\subseteq\{0,1\}^*\) as follows.  For every \(m\geq3\), its
length-\(k_m\) strings are precisely the masks
\(\boldsymbol\mu\in\{0,1\}^{k_m}\) for which the collection encoded by
\(\boldsymbol\mu\) admits an exact cover.  No strings of any other length
belong to \(L\).
An X3C collection on \([m]\) can be converted in polynomial
time to its length-\(k_m\) mask, and vice versa.  Hence \(L\) is
\(\NP\)-complete.  Since $k_m=\binom m3$ is strictly
increasing, the input length determines $m$.
For length $k_m$, use as advice
\(\pi_m:=\mathsf{Pre}(\mathbf B_{g_m})\).  The basis has
encoding length \(m^{O(1)}\), hence \(k_m^{O(1)}\).
Thus \(|\pi_m|\leq k_m^{O(1)}\), as required by
Definition~\ref{def:cvpp}.
Given $\boldsymbol\mu$, Theorem~\ref{thm:fixed-family}
computes $T_{\boldsymbol\mu}$ and $\Delta_{\boldsymbol\mu}$ in polynomial
time.  Set $\mathbf t_{\boldsymbol\mu}:=\coeff(T_{\boldsymbol\mu})$.
Definition~\ref{def:cvpp} and the correctness of the reduction give
\[
\mathsf{Dec}(\pi_m,\mathbf t_{\boldsymbol\mu},\Delta_{\boldsymbol\mu})=1
\quad\Longleftrightarrow\quad
\dist(\mathbf t_{\boldsymbol\mu},\mathcal L(I_m))^2\leq
\Delta_{\boldsymbol\mu}
\quad\Longleftrightarrow\quad
\boldsymbol\mu\in L.
\]
Therefore
$L\in\Ppoly$, implying $\NP\subseteq\Ppoly$.
\end{proof}

By the Karp--Lipton theorem~\cite{KarpLipton1980}, such a preprocessing
scheme would collapse the polynomial hierarchy to $\Sigma_2^{\mathsf P}$.

\Needspace{8\baselineskip}
\begin{remark}
Analogues of Theorem~\ref{thm:fixed-family} and
Corollary~\ref{cor:cvpp} hold after lifting the fixed family to principal
cyclic ideal lattices.  Let \(P_m,\nu_m\) be the values from
Section~\ref{sec:compiler} for \((\mathbf U_m,\mathbf 1_m)\), and set
\(H_m:=P_m^2k_m+\nu_m+4\) and
\(B_m:=2\lceil\sqrt{H_m}\rceil+1\).
We have \(\Delta_{\boldsymbol\mu}+4\leq H_m\).  Applying
Lemma~\ref{lem:principal-crt} with \(H=H_m\) and \(B=B_m\) yields one
principal cyclic ideal for each \(m\).  Its basis consists of the cyclic
shifts of its generator.  Only the lifted target and threshold depend on the
mask.  The
threshold is \(8\Delta_{\boldsymbol\mu}\), and the squared-distance gap is
\(32\).

A polynomial-time solution to exact decision-CVPP on this fixed cyclic family
would imply \(\NP\subseteq\Ppoly\).  Exact decision-CVP on the same family
remains \(\NP\)-complete, thereby answering the exact decision version of
Micciancio's question of whether CVP is \(\NP\)-hard even for a fixed family
of cyclic lattices~\cite{Micciancio2007}.
\end{remark}
